%-------------------------------------------------------------------------------
%	SECTION TITLE
%-------------------------------------------------------------------------------
\cvsection{Проекты и open-source}

%-------------------------------------------------------------------------------
%	CONTENT
%-------------------------------------------------------------------------------
\begin{cventries}

%---------------------------------------------------------
  \cventry
    {\href{https://github.com/brain-lab-research}{BRAIn Lab}} % Project
    {Администратор и контрибьютор GitHub-организации лаборатории} % Role / org
    {Git, CI/CD, GitHub Actions} % Tech
    {2024 — н.в.} % Date(s)
    {
      \begin{cvitems}
        \item {Администрирую GitHub-организацию лаборатории, выкладываю и поддерживаю код для воспроизведения экспериментов из наших статей}
      \end{cvitems}
    }

%---------------------------------------------------------
  \cventry
    {\href{https://github.com/Vepricov/LIB}{LIB}} % Project
    {Фреймворк для бенчмаркинга методов оптимизации} % Role / org
    {Python, PyTorch} % Tech
    {2024 — н.в.} % Date(s)
    {
      \begin{cvitems}
        \item {Единая платформа для сравнения оптимизаторов (AdamW, SOAP, Shampoo, Muon и др.) на задачах CV, NLP и LLM fine-tuning}
        \item {Гибкая конфигурация экспериментов, поддержка PEFT и GLUE; самый популярный из собственных репозиториев}
      \end{cvitems}
    }

%---------------------------------------------------------
  \cventry
    {\href{https://github.com/Vepricov/llm-baselines}{llm-baselines}} % Project
    {Вклад в open-source} % Role / org
    {Python} % Tech
    {2024 — н.в.} % Date(s)
    {
      \begin{cvitems}
        \item {Развиваю форк репозитория EPFL (\href{https://github.com/epfml/llm-baselines}{epfml/llm-baselines}) для предобучения (pretrain) LLM: добавляю оптимизаторы нового поколения и пайплайны экспериментов, бенчмаркаю методы на претрейне, вношу вклад в open-source}
      \end{cvitems}
    }

%---------------------------------------------------------
  \cventry
    {\href{https://github.com/Vepricov/claude-brainlab}{claude-brainlab}} % Project
    {AI-ассистент для исследовательского цикла} % Role / org
    {Claude skills, Python} % Tech
    {2025 — н.в.} % Date(s)
    {
      \begin{cvitems}
        \item {Авторская сборка из 50+ навыков и агентов для Claude Code: автоматизация всего научного цикла}
        \item {Ingest статей arXiv$\rightarrow$Zotero$\rightarrow$Obsidian, рецензирование, оформление статей и презентаций, память на графах знаний}
      \end{cvitems}
    }

%---------------------------------------------------------
  \cventry
    {\href{https://github.com/Vepricov/WLoRA}{WeightLoRA}} % Project
    {Официальный код статьи ACL 2026} % Role / org
    {Python, HuggingFace} % Tech
    {2025} % Date(s)
    {
      \begin{cvitems}
        \item {LoRA с автоматическим отбором значимых адаптеров; эксперименты на DeBERTaV3, Llama 3, Qwen (GLUE, GSM8K, SQuAD)}
        \item {Метод оформлен как PR в официальную библиотеку \href{https://github.com/huggingface/peft}{PEFT} (на ревью) — держим кулачки за одобрение}
      \end{cvitems}
    }

%---------------------------------------------------------
  \cventry
    {\href{https://github.com/Vepricov/zero-order-optimization}{zero-order-optimization}} % Project
    {Библиотека методов нулевого порядка} % Role / org
    {Python, PyTorch} % Tech
    {2025} % Date(s)
    {
      \begin{cvitems}
        \item {Набор zero-order (безградиентных) методов оптимизации в torch-совместимом формате}
      \end{cvitems}
    }

%---------------------------------------------------------
  \cventry
    {\href{https://github.com/intsystems/relaxit}{Relaxit} \& \href{https://github.com/intsystems/hippotrainer}{HippoTrainer}} % Project
    {ML-библиотеки, кафедра ИАД (МФТИ)} % Role / org
    {Python, PyTorch} % Tech
    {2025} % Date(s)
    {
      \begin{cvitems}
        \item {Relaxit — методы релаксации дискретных распределений; самый залайканный репозиторий организации кафедры ИАД (\href{https://github.com/intsystems}{intsystems}, $>290$ репозиториев), где я активный контрибьютор}
        \item {HippoTrainer — автоматический подбор гиперпараметров при обучении моделей на PyTorch}
      \end{cvitems}
    }

%---------------------------------------------------------
  \cventry
    {\href{https://github.com/Vepricov/Veprikov-MS-Thesis}{Магистерская} \& \href{https://github.com/Vepricov/Veprikov-BS-Thesis}{бакалаврская диссертации}} % Project
    {МФТИ} % Role / org
    {LaTeX} % Tech
    {2024, 2026} % Date(s)
    {
      \begin{cvitems}
        \item {Оптимальное управление петлями обратной связи в системах непрерывного обучения}
        \item {Динамическая модель эффекта обратной связи в системах ИИ}
      \end{cvitems}
    }

%---------------------------------------------------------
  \cventry
    {\href{https://github.com/Vepricov/a2pro-quant-demo}{Квантизация и ускорение инференса LLM}} % Project
    {Прикладной ML-пайплайн} % Role / org
    {Python, vLLM} % Tech
    {2025} % Date(s)
    {
      \begin{cvitems}
        \item {Квантизация Qwen3-8B (SmoothQuant + GPTQ), деплой через vLLM: $-42\%$ памяти, $+46\%$ скорости инференса без потери качества}
      \end{cvitems}
    }

\end{cventries}
